% ============================================================
%  CHAPTER 8 — Discussion
% ============================================================

\chapter{Discussion}
\label{chap:discussion}

\section{Interpretation of the principal result}

The comparison of \cref{ch:forces} separates two properties of a motor controller
that are ordinarily conflated: the \emph{pattern} of muscle recruitment and the
\emph{economy} of it. The learned controller reproduces the first and not the
second.

This is a more informative outcome than either extreme would have been. Complete
agreement would have established only that reinforcement learning can rediscover
a known optimisation criterion when the reward is arranged to favour it. Complete
disagreement would have indicated that the approach is unsuited to the problem.
The observed result --- close agreement in direction, systematic excess in
magnitude, anatomically localised to the elbow antagonists --- identifies both
what the method captures and precisely where it fails.

The localisation is itself evidence. Had the discrepancy been distributed evenly
across the nine muscles, it would have been consistent with generic noise in the
learned policy. Its concentration in two elbow extensors, driven at three to four
and a half times the prescribed force while the shoulder group agrees to within a
few per cent, identifies a specific and interpretable behaviour: sustained
co-contraction of an antagonist pair. That this is independently corroborated by
the co-contraction index (\cref{sec:icc}), computed by an unrelated method over
the same muscle groups, strengthens the inference considerably.

\section{The discount factor, and a hypothesis that had to be abandoned}

\Cref{sec:ablation2} established that the discount factor alone accounts for the
entire default-to-tuned improvement on this task: changing $\gamma$ from $0.99$ to
$0.957$ recovers $106$\,\% of the interval, while the learning rate, the
target-network rate, the batch size and the target entropy each recover nothing.

The mechanism is quantitative and checkable. The discount factor sets an
effective planning horizon of roughly $1/(1-\gamma)$ steps: 100 at the library
default, 23 at the selected value. An episode is 100 steps, but the reach itself
completes by step 25 (\cref{tab:traj}). The default therefore asked the value
function to integrate reward over four times the duration of the behaviour being
learned, and the surplus is largely holding reward --- a high-variance quantity to
estimate. A critic fed high-variance targets produces the oscillating evaluation
curves that characterised every untuned run here, and an actor trained against an
unreliable critic cannot refine its endpoint.

\subsection*{The attribution this replaces}

Earlier versions of this chapter attributed the same plateau to \SAC's default
target entropy of $-\dim(\mathcal{A})$, and argued the case at length: a
nine-muscle action space inflates the default, the retained stochasticity is
productive early and harmful when the hand must be held still, and the failure
would therefore be expected in any musculoskeletal model of comparable action
dimensionality.

The reasoning was sound and the conclusion was wrong. A two-sided ablation
(\cref{sec:ablation}) shows the target entropy recovers $-8$\,\% of the gap ---
adding it to defaults changes nothing, removing it from the tuned configuration
costs nothing. The evidence that had seemed decisive was the clustering of all
five best search trials near $-19.6$, and that clustering is precisely what a TPE
sampler produces in any region it finds promising, causal or not.

What survives is a warning of a different kind, and a more useful one. \textbf{A
hyperparameter search identifies a configuration, never a cause.} Reading
causality from which parameters cluster in the best trials is a specific,
tempting error, and separating the two requires a deliberate one-at-a-time
experiment. The diagnostic difficulty described above was real --- the symptom
does present as insufficient training or as a poorly shaped reward, and it
survived three reward redesigns, a fourfold budget correction and a doubling of
the control rate --- but the cause was the horizon, not the exploration.

\section{Tuning against algorithm selection}

\Cref{sec:fairbench} retrained all four algorithms with a correctly matched
discount factor. With every entrant configured, \SI{0.060}{\meter} separates the
best off-policy method from the second; correcting $\gamma$ alone within a single
algorithm was worth \SI{0.165}{\meter}, about three times as much.

The implication is methodological. A comparison of learning algorithms conducted
at library defaults --- the common practice, and that of the earlier version of
this work --- measures which algorithm's defaults happen to suit the problem
rather than which algorithm is better suited to it. This work demonstrated the
hazard on itself: the first benchmark ran one entrant with a corrected $\gamma$
and three without, reported \DDPG{} as clearly the weakest off-policy method, and
explained that result by its lack of twin critics. Re-run fairly, \DDPG{}
improves by $40$\,\% and becomes indistinguishable from \TDIII. The textbook
explanation had been fitted to an artefact.

That the decisive parameter here is the discount factor rather than an
algorithm-specific one is what makes the fair comparison possible at all: every
algorithm has a $\gamma$, so a common correction exists. Had the decisive
parameter been structural --- as the target entropy would have been, since only
\SAC{} has one --- there would have been no neutral setting to adopt, and the
comparison would have had no defensible form.

This does not invalidate algorithm comparison; it constrains what one can claim
from it. The defensible form tunes each algorithm's corresponding parameter under
an equal budget, which is the recommendation of \cref{ch:conclusion}.

\section{Accuracy and physiological economy as separate axes}

\Cref{sec:effortall} found that across fifteen trained policies the ordering by
muscular economy is approximately the inverse of the ordering by reaching
accuracy.

This bears on a common assumption in the application of reinforcement learning to
musculoskeletal models. Such work routinely reports task performance --- success
rate, endpoint error, completion time --- and treats good performance as evidence
that the controller behaves in a physiologically plausible manner. The present
result shows that the two properties can move in opposite directions across
algorithms trained on the same reward, on the same model, for the same task.

A controller may therefore appear excellent by every measure of what it achieves
while being unrealistic in how it achieves it, and, as \PPO demonstrates, the
converse is equally possible. If the objective is to understand or reproduce
human motor control rather than merely to displace a simulated limb, task
performance is not sufficient evidence and the muscle forces must be examined
directly.

The caveat recorded in \cref{sec:effortall} applies: \PPO's favourable economy is
not independent of its passivity, and the present data cannot fully separate the
two.

\section{Implications for the intended applications}

The clinical and assistive applications motivating this work --- prosthesis
control, exoskeleton assistance, rehabilitation devices --- require force
estimates in real time, and the original motivation for a learned controller was
that it would supply them faster than iterative solution.

\textbf{That motivation does not survive measurement.} As
\cref{sec:latencycaveat} establishes, a direct solve of the load-sharing problem
runs in roughly \SI{24}{\micro\second} against the network's
\SI{289}{\micro\second}: the classical method is about an order of magnitude
faster, and the $7.3\times$ advantage reported earlier compared the wrong
computation between two unoptimised implementations. Speed is therefore not a
reason to prefer a learned controller for this task, and the case must rest on
other grounds.

Two such grounds remain, both weaker than the original claim. The network's cost
is fixed and branch-free irrespective of biomechanical state, which matters where
a hard real-time bound is required rather than a good average. And a learned
policy maps observation directly to activation without requiring the moment-arm
matrix, the inverse-dynamics torque or a constrained solver at run time --- an
integration advantage on hardware where those are awkward to supply, though not a
computational one.

The results of \cref{ch:forces} qualify the prospect further. Where the quantity
of interest is which muscles are active and in what relative proportion --- for
instance in driving a myoelectric interface, or in classifying movement intent
--- a correlation of 0.865 may well be sufficient.
Where the absolute magnitude of individual muscle forces is the quantity of
interest --- estimating joint loading, assessing injury risk, planning surgical
intervention --- a systematic 71\,\% excess concentrated in one antagonist pair
is not acceptable, and static optimisation remains the appropriate method.

The proposed remedy of \cref{sec:whygap} is directly relevant here, since a
reward employing the classical criterion would, if the prediction holds, remove
the principal obstacle to the second class of application.

\section{Threats to validity}

The work is entirely computational. There is no physical arm, no
electromyographic recording and no human participant, and \MuJoCo's muscle
implementation is itself an approximation of the Hill formulation.

Static optimisation, the reference throughout \cref{ch:forces}, is a model of
human load sharing rather than ground truth. It encodes the assumption that the
central nervous system minimises summed squared activation, which remains
contested. A controller that departs from static optimisation is not thereby
shown to be unphysiological; it is shown to depart from the standard
computational method.

The muscle set is incomplete. The rotator cuff is absent, which required
constraining shoulder rotation (\cref{sec:rotation}), so the conclusions apply to
reaching with a restrained rotational degree of freedom.

The magnitude of the effort discrepancy is a property of the reward employed. Its
effort term is a normalised sum of squared forces weighted for numerical balance
against the distance term, which is neither the classical criterion nor scaled to
match it. That this is the operative cause is not merely argued but tested:
\cref{sec:effortsweep} reports the sweep, and raising the weight moves the effort
ratio from $1.914$ to $1.304$ against an achievable floor of $1.263$.

Finally, the statistical basis is limited throughout. The principal configuration
is replicated across five seeds and the algorithm comparison across three; no
significance test is reported anywhere in this work, because three seeds cannot
support one with useful power. \Cref{sec:accidental} gives the sharpest available
measure of the resulting uncertainty: an identical configuration, trained and
evaluated twice by this project, differed by \SI{0.036}{\meter}. These
constraints are set out in full in \cref{ch:conclusion}.
